\chapter{Stochastic Integrals with Respect to Brownian Motion}
\section{Lecture 10}

Our goal is to define the stochastic integral
\[
  \int_a^b f_t\,dB_t,
\]
where the integrand may depend on both time and the outcome $\omega$.
Because Brownian motion has paths of unbounded variation, this cannot in
general be defined as an ordinary Lebesgue--Stieltjes integral. Instead, we
first define the integral for suitably simple adapted processes and then
extend the definition by an $L^2$ limit.

\begin{definition}
Let $I=[a,b]$. Write $V(I)$ for the class of processes
$f:I\times\Omega\to\mathbb{R}$ such that:
\begin{enumerate}
  \item $f$ is $\mathcal{B}(I)\otimes\mathcal{F}$-measurable;
  \item $f$ is adapted, i.e., $f_t$ is $\mathcal{F}_t$-measurable for every
  $t\in I$;
  \item
  \[
    \mathbb{E}\left[\int_a^b f_t^2\,dt\right]<\infty.
  \]
\end{enumerate}
\end{definition}

\begin{definition}[Elementary process]
A process $\varphi\in V(I)$ is \emph{elementary} if, for a partition
\[
  a=t_0<t_1<\cdots<t_n=b,
\]
it has the form
\[
  \varphi_t=\sum_{j=0}^{n-1}A_j\mathbf{1}_{[t_j,t_{j+1})}(t),
\]
where $A_j$ is $\mathcal{F}_{t_j}$-measurable. Define
\[
  \int_a^b\varphi_t\,dB_t
  :=\sum_{j=0}^{n-1}A_j\bigl(B_{t_{j+1}}-B_{t_j}\bigr).
\]
The condition $A_j\in\mathcal{F}_{t_j}$ expresses that the value used on
the interval $[t_j,t_{j+1})$ must be chosen before its Brownian increment is
observed.
\end{definition}

\subsection*{Approximation by elementary processes}
The following three lemmas show that elementary processes are dense in $V(I)$
under the norm
\[
  \|f\|_{L^2(\mathbb{P}\times dt)}
  :=\left(\mathbb{E}\left[\int_a^b f_t^2\,dt\right]\right)^{1/2}.
\]

\begin{lemma}
Let $g\in V(I)$ be bounded and have continuous sample paths. Then there are
elementary processes $(\varphi_n)$ such that
\[
  \mathbb{E}\left[\int_a^b\bigl(g_t-\varphi_n(t)\bigr)^2\,dt\right]
  \longrightarrow0.
\]
\end{lemma}
\begin{proof}
For each $n$, choose a partition
\[
  a=t^{(n)}_0<t^{(n)}_1<\cdots<t^{(n)}_{N_n}=b
\]
whose mesh tends to zero, and set
\[
  \varphi_n(t)
  :=\sum_{j=0}^{N_n-1}g_{t_j^{(n)}}
  \mathbf{1}_{[t_j^{(n)},t_{j+1}^{(n)})}(t).
\]
Since $g$ is adapted, $g_{t_j^{(n)}}$ is
$\mathcal{F}_{t_j^{(n)}}$-measurable, so $\varphi_n$ is elementary. For each
fixed $\omega$, uniform continuity of $t\mapsto g_t(\omega)$ on $[a,b]$
implies
\[
  \int_a^b\bigl(g_t(\omega)-\varphi_n(t,\omega)\bigr)^2\,dt\to0.
\]
If $|g|\le C$, the integrands are bounded by $4C^2$. Dominated convergence
on $\Omega\times[a,b]$ proves the result.
\end{proof}

\begin{lemma}
Let $h\in V(I)$ be bounded. Then there are bounded, adapted processes
$(g_n)$ with continuous sample paths such that
\[
  \mathbb{E}\left[\int_a^b\bigl(h_t-g_n(t)\bigr)^2\,dt\right]
  \longrightarrow0.
\]
\end{lemma}
\begin{proof}
Extend $h$ by zero to times before $a$ and define the backward averages
\[
  g_n(t):=n\int_{(t-1/n)\vee a}^{t}h_s\,ds,
  \qquad a\le t\le b.
\]
Each $g_n(t)$ is $\mathcal{F}_t$-measurable, because it only uses values
$h_s$ with $s\le t$. Moreover, $g_n$ has continuous sample paths and is
bounded whenever $h$ is bounded. Translation continuity in
$L^2(\mathbb{P}\times dt)$ (equivalently, the standard convergence theorem
for one-sided moving averages) gives
\[
  g_n\longrightarrow h\quad\text{in }L^2(\mathbb{P}\times dt).
\]
This is precisely the desired assertion.
\end{proof}

\begin{lemma}
Let $f\in V(I)$. Then there are bounded processes $(h_n)\subset V(I)$ such
that
\[
  \mathbb{E}\left[\int_a^b\bigl(f_t-h_n(t)\bigr)^2\,dt\right]
  \longrightarrow0.
\]
\end{lemma}
\begin{proof}
Let
\[
  h_n(t):=(-n)\vee f_t\wedge n.
\]
The process $h_n$ remains jointly measurable and adapted, and it is bounded.
Also, $h_n(t)\to f_t$ pointwise and
\[
  |f_t-h_n(t)|^2\le f_t^2.
\]
The conclusion follows from dominated convergence, since
$f\in V(I)$.
\end{proof}

Combining the three lemmas, first truncate $f$, then smooth the bounded
truncation, and finally approximate the resulting process by elementary
processes. Hence every $f\in V(I)$ can be approximated in
$L^2(\mathbb{P}\times dt)$ by elementary processes.

\subsection*{It\^o's isometry for elementary processes}
\begin{lemma}[It\^o's isometry]
Let
\[
  \varphi_t=\sum_{j=0}^{n-1}A_j\mathbf{1}_{[t_j,t_{j+1})}(t)
\]
be elementary. Then
\[
  \mathbb{E}\left[\left(\int_a^b\varphi_t\,dB_t\right)^2\right]
  =\mathbb{E}\left[\int_a^b\varphi_t^2\,dt\right].
\]
\end{lemma}
\begin{proof}
Write $\Delta B_j=B_{t_{j+1}}-B_{t_j}$. Expanding the square gives
\[
 \mathbb{E}\left[\left(\sum_{j=0}^{n-1}A_j\Delta B_j\right)^2\right]
 =\sum_{j=0}^{n-1}\mathbb{E}[A_j^2(\Delta B_j)^2]
 +2\sum_{0\le j<k\le n-1}\mathbb{E}[A_j\Delta B_jA_k\Delta B_k].
\]
For the diagonal terms, $A_j$ is $\mathcal{F}_{t_j}$-measurable and
$\Delta B_j$ is independent of $\mathcal{F}_{t_j}$ with variance
$t_{j+1}-t_j$. Thus
\[
 \mathbb{E}[A_j^2(\Delta B_j)^2]
 =\mathbb{E}\left[A_j^2\,
 \mathbb{E}\bigl[(\Delta B_j)^2\mid\mathcal{F}_{t_j}\bigr]\right]
 =(t_{j+1}-t_j)\mathbb{E}[A_j^2].
\]
For $j<k$, the random variable $A_j\Delta B_jA_k$ is
$\mathcal{F}_{t_k}$-measurable, whereas $\Delta B_k$ is independent of
$\mathcal{F}_{t_k}$ and has conditional mean zero. Therefore
\[
 \mathbb{E}[A_j\Delta B_jA_k\Delta B_k]
 =\mathbb{E}\left[A_j\Delta B_jA_k\,
 \mathbb{E}[\Delta B_k\mid\mathcal{F}_{t_k}]\right]=0.
\]
Consequently,
\[
 \mathbb{E}\left[\left(\int_a^b\varphi_t\,dB_t\right)^2\right]
 =\sum_{j=0}^{n-1}(t_{j+1}-t_j)\mathbb{E}[A_j^2].
\]
Finally, since $\varphi_t=A_j$ on $[t_j,t_{j+1})$,
\[
 \mathbb{E}\left[\int_a^b\varphi_t^2\,dt\right]
 =\sum_{j=0}^{n-1}(t_{j+1}-t_j)\mathbb{E}[A_j^2],
\]
which proves the identity.
\end{proof}

\section{Lecture 11}
Let $V(0,T)$ denote the class of adapted, jointly measurable processes $f$ such that
\[
  \E\!\left[\int_0^T f_s^2\,ds\right]<\infty.
\]
If $(\varphi_n)$ is a sequence of elementary processes converging to $f$ in
$L^2(\Prob\times \mcal{L})$, then It\^o's isometry implies that
\[
  \left(\int_0^T\varphi_n(s)\,dB_s\right)_{n\ge1}
\]
is Cauchy in $L^2(\Prob)$. Completeness of $L^2(\Prob)$ therefore defines
$\int_0^T f_s\,dB_s$. In particular, It\^o's isometry extends from elementary
processes to every $f\in V(0,T)$:
\[
  \E\!\left[\left(\int_0^T f_s\,dB_s\right)^2\right]
  =\E\!\left[\int_0^T f_s^2\,ds\right].
\]
The stochastic integral is linear in its integrand.

\subsection*{Doob's maximal inequality}
\begin{proposition}
Let $(M_n)_{n=0}^N$ be a nonnegative submartingale. For $\lambda>0$,
\[
  \lambda\,\Prob\!\left(\max_{0\le n\le N}M_n\ge\lambda\right)
  \le \E\!\left[M_N\mathbf{1}_{\{\max_{0\le n\le N}M_n\ge\lambda\}}\right].
\]
\end{proposition}
\begin{proof}
Set $A=\{\max_{0\le n\le N}M_n\ge\lambda\}$ and let
\[
  \tau=\inf\{n\le N:M_n\ge\lambda\},
\]
with $\tau=N$ on $A^c$. By optional sampling for bounded stopping times,
\[
 \E[M_N\mid\mcal{F}_\tau]\ge M_\tau.
\]
Multiplying by $\mathbf{1}_A$ and taking expectations yields
\[
 \E[M_N\mathbf{1}_A]\ge\E[M_\tau\mathbf{1}_A]
 \ge\lambda\Prob(A).
\]
\end{proof}

\begin{corollary}
Let $(M_n)_{n=0}^N$ be a martingale or a nonnegative submartingale. For
$p\ge1$ and $\lambda>0$,
\[
 \lambda^p\Prob\!\left(\max_{0\le n\le N}|M_n|\ge\lambda\right)
 \le \E[|M_N|^p].
\]
\end{corollary}
\begin{proof}
By conditional Jensen's inequality, $(|M_n|^p)$ is a nonnegative
submartingale. Apply the proposition to this process.
\end{proof}

\begin{theorem}[Doob's maximal inequality]
Let $(M_t)_{0\le t\le T}$ be a continuous martingale. For $p\ge1$ and
$\lambda>0$,
\[
 \Prob\!\left(\sup_{0\le t\le T}|M_t|\ge\lambda\right)
 \le \frac{\E[|M_T|^p]}{\lambda^p}.
\]
\end{theorem}
\begin{proof}
Apply the discrete-time inequality to the finite sets in an increasing
sequence of finite subsets of a countable dense subset of $[0,T]$, including
$T$. Pass to the limit using continuity of $M$ and then let the usual
$\varepsilon$-approximation tend to zero.
\end{proof}

\subsection*{Continuity and martingale property}
\begin{proof}
Choose elementary processes $\varphi_n$ with
\[
 \E\!\left[\int_0^T|f_s-\varphi_n(s)|^2\,ds\right]\longrightarrow0,
 \qquad I_n(t):=\int_0^t\varphi_n(s)\,dB_s.
\]
Each $I_n$ is a continuous martingale. Hence, by Doob's inequality and
It\^o's isometry,
\[
 \Prob\!\left(\sup_{0\le t\le T}|I_n(t)-I_m(t)|>\varepsilon\right)
 \le\frac{1}{\varepsilon^2}\E\!\left[\int_0^T
 |\varphi_n(s)-\varphi_m(s)|^2\,ds\right].
\]
Choose a subsequence $(I_{n_k})$ such that
\[
 \Prob\!\left(\sup_{0\le t\le T}|I_{n_{k+1}}(t)-I_{n_k}(t)|>2^{-k}\right)<2^{-k}.
\]
Set
\[
 E_k:=\left\{\sup_{0\le t\le T}|I_{n_{k+1}}(t)-I_{n_k}(t)|>2^{-k}\right\}.
\]
Since $\sum_{k\ge1}\Prob(E_k)\le\sum_{k\ge1}2^{-k}<\infty$, the
Borel--Cantelli lemma implies that, almost surely, only finitely many of the
events $E_k$ occur. Hence, for almost every $\omega$, there is an integer
$K(\omega)$ such that, for every $k\ge K(\omega)$,
\[
 \sup_{0\le t\le T}|I_{n_{k+1}}(t,\omega)-I_{n_k}(t,\omega)|\le2^{-k}.
\]
For $m>\ell\ge K(\omega)$, the triangle inequality now yields
\[
 \begin{aligned}
 \sup_{0\le t\le T}|I_{n_m}(t,\omega)-I_{n_\ell}(t,\omega)|
 &\le\sum_{k=\ell}^{m-1}\sup_{0\le t\le T}
 |I_{n_{k+1}}(t,\omega)-I_{n_k}(t,\omega)|\\
 &\le\sum_{k=\ell}^{m-1}2^{-k}.
 \end{aligned}
\]
The final quantity tends to zero as $\ell\to\infty$. Thus, for almost every
$\omega$, the continuous functions $I_{n_k}(\cdot,\omega)$ are Cauchy in
the supremum norm on $[0,T]$. They therefore converge uniformly to a
continuous function $I(\cdot,\omega)$. Define $I$ arbitrarily on the null
set where this argument fails.

For every fixed $t$, $I_{n_k}(t)\to I_t$ almost surely, while
$I_{n_k}(t)\to M_t$ in $L^2$, hence in probability. Uniqueness of limits in
probability gives $I_t=M_t$ almost surely. Thus $I$ is a continuous version
of $M$.

For $s\le t$, conditional expectation is continuous in $L^2$, and therefore
\[
 \E[M_t\mid\F_s]
 =\lim_{k\to\infty}\E[I_{n_k}(t)\mid\F_s]
 =\lim_{k\to\infty}I_{n_k}(s)=M_s.
\]
Finally, apply Doob's inequality with $p=2$ and use It\^o's isometry.
\end{proof}

\section{Lecture 12}
\subsection*{Stopping an It\^o integral}
\begin{lemma}
Let $f\in V(0,T)$ and let $\tau$ be a stopping time with values in $[0,T]$.
If $M_t=\int_0^t f_s\,dB_s$, then
\[
 M_\tau=\int_0^\tau f_s\,dB_s
 =\int_0^T f_s\mathbf{1}_{\{s<\tau\}}\,dB_s.
\]
\end{lemma}
\begin{proof}
The identity holds directly when $f$ is elementary and $\tau$ is a simple
stopping time. Let $f_n$ be elementary processes with
\[
 \E\!\left[\int_0^T|f_n(s)-f_s|^2\,ds\right]\to0,
\]
and let simple stopping times $\tau_n$ converge pointwise to $\tau$. Then
\[
 \int_0^{\tau_n}f_n(s)\,dB_s
 =\int_0^T f_n(s)\mathbf{1}_{\{s<\tau_n\}}\,dB_s.
\]
The left-hand integrands converge in $L^2(\Prob\times\mcal{L})$ to
$f(s)\mathbf{1}_{\{s<\tau\}}$: one term follows from the approximation
$f_n\to f$, and the other from dominated convergence, since
$\mathbf{1}_{\{s<\tau_n\}}\to\mathbf{1}_{\{s<\tau\}}$ for Lebesgue-a.e.
$s$. It\^o's isometry consequently gives convergence in $L^2$ of the
right-hand side.

On the other hand, Doob's inequality gives uniform convergence in probability
of $\int_0^\cdot f_n\,dB$ to $\int_0^\cdot f\,dB$. Together with continuity
of the latter process and $\tau_n\to\tau$, this yields
\[
 \int_0^{\tau_n} f_n(s)\,dB_s\ \xrightarrow{\Prob}\ \int_0^\tau f_s\,dB_s.
\]
The two limits identify the asserted identity.
\end{proof}

\subsection*{Multidimensional It\^o integrals}
Let $B=(B^1,\ldots,B^n)$ be an $n$-dimensional Brownian motion; its
coordinates are independent one-dimensional Brownian motions. Let
$V^{m\times n}(a,b)$ be the class of $m\times n$ matrix-valued processes
$V=(V_{ij})$ for which every entry is jointly measurable, adapted, and
\[
 \E\!\left[\int_a^b V_{ij}(t)^2\,dt\right]<\infty.
\]
For $V\in V^{m\times n}(a,b)$, define
\[
 \int_a^bV_t\,dB_t
 :=\begin{pmatrix}
 \sum_{j=1}^n\int_a^bV_{1j}(t)\,dB_t^j\\
 \vdots\\
 \sum_{j=1}^n\int_a^bV_{mj}(t)\,dB_t^j
 \end{pmatrix}.
\]
Thus the integral is an $m\times1$ random vector. Equivalently, it is the
matrix product $\int_a^b V_t\,dB_t$ interpreted componentwise.

\subsection*{Locally square-integrable integrands}
The global condition
\[
 \E\!\left[\int_a^b f_t^2\,dt\right]<\infty
\]
can be weakened to
\[
 \Prob\!\left(\int_a^b f_t^2\,dt<\infty\right)=1.
\]
Write $\mathcal W(a,b)$ for the class of jointly measurable, adapted
processes satisfying this almost-sure condition; similarly,
$\mathcal W^{m\times n}(a,b)$ denotes the matrix-valued version.

One construction uses the stopping times
\[
 \tau_n:=\inf\!\left\{t\ge a:\int_a^t f_s^2\,ds\ge n\right\}\wedge b,
 \qquad f^{(n)}_t:=f_t\mathbf{1}_{\{t<\tau_n\}}.
\]
Then $\tau_n\uparrow b$ almost surely and each $f^{(n)}$ belongs to
$V(a,b)$. Define
\[
 \int_a^b f_t\,dB_t
 :=\lim_{n\to\infty}\int_a^b f^{(n)}_t\,dB_t.
\]
The limit is well defined because the stopped integrals agree up to
$\tau_n$ whenever the later stopping time is used.

\begin{remark}
The usual algebraic and stopping-time properties of the stochastic integral
extend to $\mathcal W(a,b)$. The process
$M_t=\int_a^t f_s\,dB_s$ need not be a true martingale, because the
construction only gives convergence in probability rather than global
$L^2$ convergence. It is, however, a local martingale: there are stopping
times $\tau_n\uparrow\infty$ almost surely for which
$(M_{t\wedge\tau_n})_{t\ge0}$ is a martingale for every $n$.
\end{remark}
\section{Lecture 13}

Let $(B_t)_{t\ge0}$ be a one-dimensional Brownian motion on
$(\Omega,\F,\Prob)$.

\begin{definition}[It\^o process]
A one-dimensional It\^o process is a stochastic process of the form
\[
 X_t=X_0+\int_0^t u_s\,ds+\int_0^t v_s\,dB_s,
\]
where $v$ is locally square-integrable and $u$ is adapted with
\[
 \int_0^t|u_s|\,ds<\infty\qquad\text{almost surely for every }t\ge0.
\]
In differential notation, we write
\[
 dX_t=u_t\,dt+v_t\,dB_t.
\]
The process $u$ is called the \emph{drift} and $v$ the \emph{diffusion
coefficient} (or volatility).
\end{definition}

If $X$ is an It\^o process, then, whenever the two terms are well defined,
the stochastic integral with respect to $X$ is defined by
\[
 \int_0^t f_s\,dX_s
 :=\int_0^t f_su_s\,ds+\int_0^t f_sv_s\,dB_s.
\]
For instance, this holds if $f$ is adapted and almost surely bounded on each
finite time interval.

\begin{theorem}[One-dimensional It\^o formula]
Let $X$ satisfy $dX_t=u_t\,dt+v_t\,dB_t$, and let
$g\in C^{1,2}([0,\infty)\times\R)$. Then
$Y_t:=g(t,X_t)$ is an It\^o process and
\[
 \begin{aligned}
 g(t,X_t)=g(0,X_0)
 &+\int_0^t g_x(s,X_s)v_s\,dB_s\\
 &+\int_0^t\left[g_t(s,X_s)+g_x(s,X_s)u_s
 +\frac12g_{xx}(s,X_s)v_s^2\right]ds.
 \end{aligned}
\]
\end{theorem}

\begin{remark}[Differential notation]
The formula is commonly written as
\[
 dY_t=g_x(t,X_t)v_t\,dB_t+
 \left[g_t(t,X_t)+g_x(t,X_t)u_t
 +\frac12g_{xx}(t,X_t)v_t^2\right]dt.
\]
Equivalently,
\[
 dY_t=g_t(t,X_t)\,dt+g_x(t,X_t)\,dX_t
 +\frac12g_{xx}(t,X_t)(dX_t)^2,
\]
using the mnemonic rules
\[
 (dt)^2=dt\,dB_t=dB_t\,dt=0,\qquad(dB_t)^2=dt.
\]
Thus $(dX_t)^2=v_t^2dt$. These are shorthand for limits of sums of
increments, not ordinary algebraic identities.
\end{remark}

In the simplest case, $Y_t=g(B_t)$ with $g\in C^2(\R)$, It\^o's formula
becomes
\[
 dY_t=g'(B_t)\,dB_t+\frac12g''(B_t)\,dt.
\]
The proof was deferred to the more general theory of continuous
semimartingales.
\section{Lecture 14}

\begin{example}[Geometric Brownian motion]
Let
\[
 X_t=X_0\exp\!\left(\left(r-\frac12\sigma^2\right)t+\sigma B_t\right).
\]
Set $g(t,x)=\exp((r-\tfrac12\sigma^2)t+\sigma x)$. Applying It\^o's
formula to $X_t/X_0=g(t,B_t)$ gives
\[
 \begin{aligned}
 d\left(\frac{X_t}{X_0}\right)
 &=\sigma g(t,B_t)\,dB_t
 +\left(r-\frac12\sigma^2\right)g(t,B_t)\,dt
 +\frac12\sigma^2g(t,B_t)\,dt\\
 &=\sigma\frac{X_t}{X_0}\,dB_t+r\frac{X_t}{X_0}\,dt.
 \end{aligned}
\]
Consequently,
\[
 X_t=X_0+\int_0^t\sigma X_s\,dB_s+\int_0^t rX_s\,ds,
\]
or, in differential notation,
\[
 dX_t=rX_t\,dt+\sigma X_t\,dB_t.
\]
\end{example}

\subsection*{Integration by parts}
\begin{theorem}[Stochastic product rule]
Let $X$ and $Y$ be It\^o processes driven by the same Brownian motion:
\[
 dX_t=u_1(t)\,dt+v_1(t)\,dB_t,
 \qquad dY_t=u_2(t)\,dt+v_2(t)\,dB_t.
\]
Then
\[
 d(X_tY_t)=X_t\,dY_t+Y_t\,dX_t+v_1(t)v_2(t)\,dt.
\]
Equivalently,
\[
 \int_0^tX_s\,dY_s
 =X_tY_t-X_0Y_0-\int_0^tY_s\,dX_s
 -\int_0^t v_1(s)v_2(s)\,ds.
\]
\end{theorem}
The final integral is the difference from the classical integration-by-parts
formula. It arises because $dB_t\,dB_t=dt$.

\subsection*{Quadratic variation}
For Brownian motion and partitions $\Pi_n$ of $[a,b]$ whose mesh tends to
zero,
\[
 \sum_{[t_i,t_{i+1}]\in\Pi_n}(B_{t_{i+1}}-B_{t_i})^2
 \xrightarrow{L^2} b-a.
\]
This explains the heuristic $(dB_t)^2=dt$. For example, It\^o's formula
gives
\[
 d(B_t^2)=2B_t\,dB_t+dt,
\]
so $B_t^2-t$ has no drift term and is a martingale.

\begin{theorem}[Quadratic variation of a continuous martingale]
Let $M$ be a continuous $L^2$-bounded martingale with $M_0=0$ almost surely.
Then there exists a unique (up to indistinguishability) continuous,
increasing adapted process $[M]$, with $[M]_0=0$, such that
\[
 M_t^2-[M]_t
\]
is a uniformly integrable martingale. The process $[M]$ is called the
\emph{quadratic variation} of $M$.
\end{theorem}

For a partition $\Delta=\{t_0<t_1<\cdots\}$ having finitely many points in
each bounded interval, define the realized squared-increment process
\[
 T_t^\Delta(M)
 :=\sum_{i=0}^{k-1}(M_{t_{i+1}}-M_{t_i})^2+(M_t-M_{t_k})^2,
\]
where $t_k\le t<t_{k+1}$. The omitted proof establishes the appropriate
limit of $T_t^\Delta(M)$ as the mesh tends to zero and identifies it with
$[M]_t$.
\section{Lecture 15: Quadratic Variation and Local Martingales}

Throughout this section, let \(M=\{M_t\}_{t\ge 0}\) be a continuous, \(L^2\)-bounded martingale with \(M_0=0\). We first treat the case in which \(M\) is bounded.

Fix \(a>0\), and let
\[
\Delta=\{0=t_0<t_1<\cdots<t_k=a\}
\]
be a partition of \([0,a]\). For \(t\in[t_j,t_{j+1})\), define
\[
T_t^\Delta(M)
=
\sum_{i=0}^{j-1}(M_{t_{i+1}}-M_{t_i})^2
+
(M_t-M_{t_j})^2.
\]
Thus, at a partition point \(t_j\),
\[
T_{t_j}^\Delta(M)
=
\sum_{i=0}^{j-1}(M_{t_{i+1}}-M_{t_i})^2.
\]

For each fixed partition \(\Delta\), the process
\[
\{M_t^2-T_t^\Delta(M)\}_{0\le t\le a}
\]
is a martingale. Indeed, on an interval \([t_j,t_{j+1})\),
\[
M_t^2-T_t^\Delta(M)
=
\left(M_{t_j}^2-\sum_{i=0}^{j-1}(M_{t_{i+1}}-M_{t_i})^2\right)
+
2M_{t_j}(M_t-M_{t_j}),
\]
and the final term is a martingale on that interval.

Let \(\{\Delta_n\}\) be a sequence of partitions of \([0,a]\) whose meshes tend to zero. Earlier, we showed that
\[
\{T_a^{\Delta_n}(M)\}_{n\ge 1}
\]
is Cauchy in \(L^2\). Since
\[
T^{\Delta_n}(M)-T^{\Delta_m}(M)
\]
is a martingale, Doob's \(L^2\) maximal inequality gives
\[
\mathbb P\left(
\sup_{0\le s\le a}
\left|T_s^{\Delta_n}(M)-T_s^{\Delta_m}(M)\right|
\ge \lambda
\right)
\le
\frac{1}{\lambda^2}
\mathbb E\left[
\left(T_a^{\Delta_n}(M)-T_a^{\Delta_m}(M)\right)^2
\right].
\]
Consequently, \(\{T^{\Delta_n}(M)\}\) is Cauchy in probability for the supremum norm on \([0,a]\).

We may therefore select a subsequence \(\{\Delta_{n_k}\}\) such that
\[
\mathbb P\left(
\sup_{0\le s\le a}
\left|T_s^{\Delta_{n_{k+1}}}(M)-T_s^{\Delta_{n_k}}(M)\right|
\ge 2^{-k}
\right)
\le 2^{-k}.
\]
Let
\[
A_k=
\left\{
\sup_{0\le s\le a}
\left|T_s^{\Delta_{n_{k+1}}}(M)-T_s^{\Delta_{n_k}}(M)\right|
\ge 2^{-k}
\right\}.
\]
Since
\[
\sum_{k=1}^{\infty}\mathbb P(A_k)<\infty,
\]
the Borel--Cantelli lemma implies that, almost surely, only finitely many \(A_k\) occur. Hence, almost surely,
\[
\{T^{\Delta_{n_k}}(M)\}_{k\ge 1}
\]
is Cauchy in the supremum norm on \([0,a]\). It therefore converges uniformly on \([0,a]\) to a continuous process, denoted by \([M]\).

The limit does not depend on the particular sequence of partitions with mesh tending to zero. Repeating this construction on \([0,1],[0,2],\ldots\) and using a diagonal subsequence gives a continuous process \([M]=\{[M]_t\}_{t\ge0}\).

We now show that \([M]\) is increasing. Let \(0\le s<t\). Since the limit is independent of the choice of partitions, we may choose partitions containing both \(s\) and \(t\). Then
\[
T_s^\Delta(M)\le T_t^\Delta(M),
\]
because at partition points these quantities are partial sums of nonnegative squared increments. Taking limits yields
\[
[M]_s\le [M]_t.
\]
Also, \([M]_0=0\).

For bounded \(M\), we may pass to the limit in
\[
M_t^2-T_t^\Delta(M).
\]
Thus,
\[
M_t^2-[M]_t
\]
is a martingale. Moreover,
\[
\mathbb E([M]_t)=\mathbb E(M_t^2).
\]
Since \([M]\) is increasing and \(M\) is bounded,
\[
\mathbb E([M]_\infty)
=
\lim_{t\to\infty}\mathbb E([M]_t)
=
\lim_{t\to\infty}\mathbb E(M_t^2)
<\infty.
\]
Hence \([M]_t\le [M]_\infty\) for every \(t\), and the family
\[
\{M_t^2-[M]_t:t\ge0\}
\]
is uniformly integrable.

\begin{theorem}[Uniqueness of quadratic variation]
Suppose that \(Y^1\) and \(Y^2\) are continuous increasing processes, both vanishing at zero, such that
\[
M^2-Y^1
\qquad\text{and}\qquad
M^2-Y^2
\]
are martingales. Then \(Y^1=Y^2\) almost surely.
\end{theorem}

\begin{proof}
Set
\[
Z=Y^1-Y^2.
\]
Then
\[
Z=(M^2-Y^2)-(M^2-Y^1),
\]
so \(Z\) is a continuous martingale with \(Z_0=0\). Since \(Y^1\) and \(Y^2\) both have finite variation, \(Z\) has finite variation. Therefore,
\[
[Z]=0.
\]
Since \(Z^2-[Z]=Z^2\) is a martingale,
\[
\mathbb E(Z_t^2)=\mathbb E(Z_0^2)=0.
\]
Thus \(Z_t=0\) almost surely for every \(t\), and hence \(Y^1=Y^2\).
\end{proof}

We now remove the assumption that \(M\) is bounded. Define
\[
T_n=\inf\{t\ge0:|M_t|\ge n\}.
\]
Since \(M\) is continuous, the stopped martingale
\[
M^{T_n}_t=M_{t\wedge T_n}
\]
is bounded by \(n\). Therefore it has a quadratic variation process, denoted by \(A^n\), such that
\[
(M^{T_n})^2-A^n
\]
is a martingale.

For \(n<m\), stopping the martingale
\[
(M^{T_m})^2-A^m
\]
at time \(T_n\) gives
\[
(M^{T_n})^2-(A^m)^{T_n},
\]
which is also a martingale. By uniqueness for bounded martingales,
\[
A^n=(A^m)^{T_n}.
\]
Thus these processes agree consistently before the corresponding stopping times. We may define
\[
[M]_t=A^n_t
\qquad\text{on } \{t\le T_n\}.
\]
This definition is independent of \(n\), and \([M]\) is continuous, increasing, and satisfies \([M]_0=0\).

\begin{definition}[Local martingale]
An adapted process \(X=\{X_t\}_{t\ge0}\) is called an
\((\mathcal F_t,\mathbb P)\)-local martingale if there exist stopping times
\[
T_n\uparrow\infty
\qquad\text{almost surely}
\]
such that \(X^{T_n}\) is a martingale for every \(n\). The sequence
\(\{T_n\}\) is called a \emph{localizing sequence} or \emph{reducing sequence}.
\end{definition}

A standard example of a strict local martingale is obtained from three-dimensional Brownian motion \(B\), started away from the origin:
\[
X_t=\frac{1}{\|B_t\|}.
\]
This is a local martingale. For example, one may localize it using
\[
T_n=\inf\{t\ge0:X_t\ge n\}.
\]
\section*{Lecture 16: Class \texorpdfstring{$DL$}{DL}, Quadratic Covariation, and Elementary It\^o Integrals}

\subsection*{When is a local martingale a martingale?}

We have already encountered many local martingales of the form
\[
    \int_0^t H_s\,dB_s.
\]
Under a suitable integrability condition, such a process is a true martingale.
The following criterion gives a useful general answer.

\begin{definition}[Class \(DL\)]
An adapted process \(X=\{X_t\}_{t\geq0}\) is in \emph{class \(DL\)} if, for every \(a>0\), the family
\[
    \{X_T:T\text{ is a stopping time and }0\leq T\leq a\}
\]
is uniformly integrable.
\end{definition}

The notation \(DL\) is commonly read as either \emph{Doob localized} or
\emph{dominated localized}.

\begin{lemma}
Let \(Y\in L^1\), and let \(\{\mathcal G_\alpha\}_{\alpha\in I}\) be any collection of sub-\(\sigma\)-algebras of \(\mathcal F\). Then
\[
    \left\{\E[Y\mid\mathcal G_\alpha]:\alpha\in I\right\}
\]
is uniformly integrable.
\end{lemma}

\begin{theorem}
A continuous local martingale \(X\) is a martingale if and only if it is in class \(DL\).
\end{theorem}

\begin{proof}
Suppose first that \(X\) is a martingale, and let \(T\) be a stopping time
with \(0\leq T\leq a\). Optional sampling gives
\[
    X_T=\E[X_a\mid\F_T].
\]
The lemma therefore shows that \(\{X_T:0\leq T\leq a\}\) is uniformly
integrable. Since \(a>0\) was arbitrary, \(X\) is in class \(DL\).

Conversely, let \(X\) be a local martingale in class \(DL\), and let
\(\{T_n\}\) be a localizing sequence. Fix \(0\leq s\leq t\) and
\(A\in\F_s\). Since \(X^{T_n}\) is a martingale,
\[
    \E[X_{t\wedge T_n};A]=\E[X_{s\wedge T_n};A].
\]
Because \(T_n\uparrow\infty\) almost surely and \(X\) has continuous paths,
\[
    X_{t\wedge T_n}\longrightarrow X_t,
    \qquad
    X_{s\wedge T_n}\longrightarrow X_s
\]
almost surely. Both sequences are uniformly integrable by the \(DL\)
assumption, so the convergences also hold in \(L^1\). Hence
\[
    \E[X_t;A]=\E[X_s;A].
\]
As this holds for every \(A\in\F_s\), we have
\(\E[X_t\mid\F_s]=X_s\) almost surely. Thus \(X\) is a martingale.
\end{proof}

\subsection*{The process \texorpdfstring{$M^2-[M]$}{M\textsuperscript{2} - [M]} for an \texorpdfstring{$L^2$}{L2}-bounded martingale}

Let \(M\) be a continuous, \(L^2\)-bounded martingale with \(M_0=0\).
We show that \(M^2-[M]\) is a uniformly integrable martingale.

Since \(M\) is \(L^2\)-bounded, there exists \(M_\infty\in L^2\) such that
\(M_t\to M_\infty\) almost surely and in \(L^2\). Doob's \(L^2\) maximal
inequality yields
\[
    \E\bigl[(M_\infty^*)^2\bigr]
    \leq 4\E[M_\infty^2]<\infty,
    \qquad
    M_\infty^*:=\sup_{t\geq0}|M_t|.
\]

For \(n\in\mathbb N\), define
\[
    T_n:=\inf\{t\geq0:|M_t|\geq n\}.
\]
The stopped martingale \(M^{T_n}\) is bounded. Therefore, for each \(t\geq0\),
\[
    \E\bigl[[M]_{t\wedge T_n}\bigr]
    =\E\bigl[M_{t\wedge T_n}^2\bigr]
    \leq\E\bigl[(M_\infty^*)^2\bigr].
\]
Letting \(n\to\infty\), monotone convergence gives
\[
    \E[[M]_t]\leq\E\bigl[(M_\infty^*)^2\bigr].
\]
One more application of monotone convergence, now as \(t\to\infty\), gives
\[
    \E[[M]_\infty]
    \leq\E\bigl[(M_\infty^*)^2\bigr]<\infty.
\]
Consequently,
\[
    \sup_{t\geq0}|M_t^2-[M]_t|
    \leq (M_\infty^*)^2+[M]_\infty
\]
is integrable. Thus \(M^2-[M]\) is in class \(DL\). Since it is already a
continuous local martingale, the preceding theorem implies that it is a
martingale; its class \(DL\) property also makes this martingale uniformly
integrable.

\subsection*{Quadratic covariation}

\begin{definition}[Quadratic covariation]
Let \(M\) and \(N\) be continuous, \(L^2\)-bounded martingales with
\(M_0=N_0=0\). Their \emph{quadratic covariation} is defined by polarization:
\[
    [M,N]:=\frac14\bigl([M+N]-[M-N]\bigr).
\]
\end{definition}

In particular, \([M,M]=[M]\).

\begin{theorem}
The process \([M,N]\) is the unique continuous finite-variation process
vanishing at zero such that
\[
    MN-[M,N]
\]
is a uniformly integrable martingale.
\end{theorem}

\begin{proof}
Omitted.
\end{proof}
